% Introduction

Reinforcement learning for control aims to learn policies through interaction with dynamic environments. In this project, the long-term objective is to investigate whether Group Relative Policy Optimization (GRPO), originally introduced in the context of language-model reinforcement learning and mathematical reasoning, can be adapted to control-oriented benchmark environments \citep{shao2024deepseekmath}. Such an adaptation is non-trivial because control tasks require sequential credit assignment, stable learning from trajectories, and evaluation across environments, random seeds, and training budgets.

The final project will compare a GRPO-based control method against PPO, SAC, and TD3. These algorithms represent widely used baselines in deep reinforcement learning: PPO is an on-policy policy-gradient method, SAC is an off-policy maximum-entropy actor-critic method, and TD3 is an off-policy deterministic actor-critic method for continuous control \citep{schulman2017proximal,haarnoja2018sacapps,fujimoto2018td3}. In this midway report, however, the GRPO-control method has not yet been implemented. The contribution is therefore to establish the formal notation, implementation pipeline, and baseline result matrix required for the final comparison.

The midway implementation completes PPO, SAC, and TD3 baselines on three environments: \texttt{cartpole\_swingup}, \texttt{acrobot\_swingup}, and \texttt{car\_racing\_continuous}. The two vector-control environments are run through ObjectRL, while CarRacing is handled by project-side CNN implementations because the image-observation setup does not fit directly into the standard vector-observation workflow \citep{baykal2025objectrl,towers2024gymnasium}. The completed baseline matrix contains three algorithms, three environments, and five seeds, corresponding to 45 runs and 900 evaluation rows.

The results should be interpreted as an interim foundation rather than final performance claims. No hyperparameter optimization was performed, and the training budgets were chosen to validate the pipeline rather than to establish converged performance. The main outcome of this report is therefore a reproducible experimental foundation: the MDP notation is fixed, the baseline matrix is complete, the result files are validated, and learning curves can be generated from the stored evaluation data.

The remainder of the report is organized as follows. Section~\ref{sec:methodology} introduces the MDP notation and experimental protocol. Section~\ref{sec:experiments} presents the completed midway baseline experiments. The final report will extend the related work, theory, and conclusion sections when the GRPO-control method has been formalized and evaluated.
